\documentstyle[12pt]{article}

%\renewcommand{\baselinestretch}{1.25}



\setlength{\headheight}{0in}
\setlength{\headsep}{0in}
\setlength{\oddsidemargin}{0.35in}
\setlength{\topmargin}{0.25in}
\setlength{\textheight}{8.5in}
\setlength{\textwidth}{5.75in}



\begin{document}


\centerline{\bf CURRICULUM VITAE}

\vspace{2 cm}

\noindent\hbox to 2.5in{\bf Nom:\hfill} ALAIN GUAY\\
%\vspace{.25 cm}

\noindent\hbox to 2.5in{\bf Adresse au travail:\hfill} Universit\'e du Qu\'ebec \`a Montr\'eal\\
\noindent\hbox to 2.5in{\hfill} D\'epartement des sciences \'economiques\\
\noindent\hbox to 2.5in{\hfill} Case postale 8888, Succursale Centre-Ville\\
\noindent\hbox to 2.5in{\hfill} Montr\'eal (Qc), Canada, H3C 3P8\\
%\vspace{.25 cm}

\noindent\hbox to2.5in{\bf T\'el\'ephone:\hfill} (Bureau) (514) 987-3000 \# 8377\\
\noindent\hbox to2.5in{\hfill} (Fax) (514) 987-8494\\
\noindent\hbox to 2.5in{\hfill} (e-mail) guay.alain@uqam.ca\\
%\vspace{.25 cm}

\noindent\hbox to 2.5in{\bf Citoyennet\'e:\hfill} Canadienne\\
%\vspace{.5 cm}

\noindent\hbox to 2.5in{\bf Langues:\hfill} Fran\c cais, anglais
%\end{titlepage}
\vspace{1 cm}

\centerline{\bf Formation:}
\vspace{.5 cm}

\noindent\hbox to 1.3in{\bf1989-1994 :\hfill} Doctorat en sciences \'economiques, Universit\'e de Montr\'eal
%\vspace{.5 cm}

\noindent\hbox to1.3in{\bf 1988 :\hfill} Certificat en statistique, Universit\'e Laval
%\vspace{.5 cm}

\noindent\hbox to1.3in{\bf1984-1986 :\hfill} M.Sc. en \'economie, Universit\'e de Montr\'eal
%\vspace{.5 cm}

\noindent\hbox to1.3in{\bf1981-1984 :\hfill} Baccalaur\'eat en \'economie, Universit\'e Laval.

\vspace{.5 cm}



\centerline{\bf Exp\'eriences professionnelles:}
\vspace{.5 cm}

\noindent\hbox to 1.3in{\bf 2019     :\hfill} Cotitulaire de la Chaire en macro\'economie et pr\'evisions,\\
\noindent\hbox to 1.3in{\hfill} ESG-UQAM

\noindent\hbox to 1.3in{\bf 2004     :\hfill} Professeur titulaire,\\
\noindent\hbox to 1.3in{\hfill} Universit\'e du Qu\'ebec \`a Montr\'eal


\noindent\hbox to 1.3in{\bf 1999-2004:\hfill} Professeur aggr\'eg\'e,\\
\noindent\hbox to 1.3in{\hfill} Universit\'e du Qu\'ebec \`a Montr\'eal

\noindent\hbox to 1.3in{\bf 1996-1999:\hfill} Professeur associ\'e,\\
\noindent\hbox to 1.3in{\hfill} Universit\'e du Qu\'ebec \`a Montr\'eal

%\vspace{.5cm}

\noindent\hbox to 1.3in{\bf 1995:\hfill} \'Economiste, Minist\`ere des finances, Canada\\
\noindent\hbox to 1.3in{\hfill} D\'epartement des \'etudes sp\'eciales

%\vspace{.5cm}

\noindent\hbox to 1.3in{\bf 1994:\hfill} \'Economiste, Banque du Canada\\
\noindent\hbox to 1.3in{\hfill} D\'epartement des relations internationales

%\vspace{.5cm}



%\vfill\eject
\newpage

\vspace{.3 cm}

\centerline{\bf Chercheur et professeur invit\'es:}

\vspace{.3 cm}

\begin{description}
\item Professeur invit\'e: Universit\'e Paris Ouest, Nanterre: novembre 2013, octobre 2014, d\'ecembre 2015, novembre 2016, octobre 2017, octobre 2018, octobre 2019, octobre 2021, octobre 2022, octobre 2023, octobre 2024, octobre 2025 et octobre 2026.
\item Chercheur invit\'e: Toulouse School of Economics (TSE), Universit\'e de Toulouse : janvier \`a
juillet 2003, octobre 2004, octobre 2005, octobre 2007, mai 2008, mai 2009, novembre 2010, mai 2012, juin 2013, juin 2014, juin 2015, juin 2016, juin 2017, mai 2018, juin 2019, juin 2022, juin 2023, juin 2024, juin 2025, mai 2026 et juin 2026.
\item Chercheur invit\'e: Centre de Recerca en Economia Internacional (CREI), Universitat Pompeu Fabra : septembre 2010 \`a mai 2011.
\item Professeur invit\'e: Paris-I, Panth\'eon-Sorbonne : mai 2006.
\item Professeur invit\'e: HEC Lausanne : mars 2006, mai 2008 et mai 2009.
\item Professeur invit\'e: Universit\'e de la R\'eunion : mars 2008, novembre 2008, mars 2009, mars 2010, mars 2011, mars 2012, mars 2013 et mars 2014.
\item Professeur invit\'e: Universit\'e des Antilles : mars 2004, mars 2005, f\'evrier 2006, d\'ecembre 2006,
janvier 2007, novembre 2007, janvier 2008, janvier 2009, janvier 2010, janvier 2011, d\'ecembre 2011, d\'ecembre 2012, septembre 2012, d\'ecembre 2013, mai 2014, d\'ecembre 2014, f\'evrier 2015, mai 2016, d\'ecembre 2016, mars 2017, d\'ecembre 2017, mars 2018, d\'ecembre 2019, d\'ecembre 2021, janvier 2023, d\'ecembre 2023, d\'ecembre 2024, d\'ecembre 2025 et d\'ecembre 2026.
\end{description}


\vspace{.3 cm}

\newpage

\centerline{\bf Publications:}
\vspace{.3 cm}


\begin{description}

\item[]  Bec, F., A. Guay, H. B. Nielsen et S. Sa\"{\i}di (2025) ``Power of unit root tests against
nonlinear and noncausal alternatives,
with an application to the Brent crude oil price'', {\it Studies in Nonlinear Dynamics \& Econometrics,} 29 (1), 1--18.


\item[] Guay, A. et F. Pelgrin (2023), ``Structural VAR Models in the Frequency Domain'', {\it Journal of Econometrics}, 236, article 105466.


\item[] Guay, A. (2021), ``Identification of Structural Vector Autoregressions Through Higher Unconditional Moments", {\it Journal of Econometrics}, 225, 27--46.

\item[] Beaudry, P., P. F\`eve, A. Guay et F. Portier (2019), ``When is Nonfundamentalness in SVARs a Real Problem?", {\it Review of Economic Dynamics,} 30, 221--243.


\item[]  F\`eve, P. et A. Guay (2019), ``Sentiments in SVARs", {\it Economic Journal,} 129, 877--896.


%\item F\`eve, P. et A. Guay (2018), ``Sentiments in SVARs", {\sl Economic Journal,} \\ https://doi.org/10.1111/ecoj.12580.

\item[] Guay, A. et J.F. Lamarche (2018), ``Structural Change Tests for GEL Criteria", {\it Econometric Reviews,}
37, 1000--1032.

\item[] Guay, A. et F. Pelgrin (2016), ``Using implied probabilities to improve estimation with
unconditional moment restrictions for weakly dependent data",  {\it Econometric Reviews,} 35, 344--372.

\item[] Guay, A. (2016), ``Anticipations, bruits et sentiments,"   {\it L'Actualit\'e \'Economique: revue d'analyse \'economique,}  92, 621--648.

\item[] Guay, A. et A. Maurin (2015), ``Disaggregation Methods Based on MIDAS Regression'', {\it Economic Modelling,} 50, 123--172.

\item[] Chaudourne, J., P. F\`eve et A. Guay (2014), ``Understanding the Effect of Technology Shocks in SVARs with
Long-Run Restrictions'', {\it Journal of Economic Dynamics and Control,} 41, 154--172.

\item[] Guay, A., E. Guerre et S. Lazarova, (2013) ``Robust Adaptive Rate-Optimal Testing for the White Noise  
Hypothesis", {\it Journal of Econometrics,} 176, 134--145.

\item[] Ambler, S., A. Guay et L. Phaneuf, (2012), ``Endogenous Business-Cycle Propagation and the Persistence Problem: The Role of Labor-Market Frictions'',   {\it Journal of Economic Dynamics and Control}, 36, 47--62.

\item[] Guay, A. et J.F. Lamarche (2012), ``Structural Change Tests Based on 
Implied Probabilities for GEL Criteria",  {\it Econometric Theory}, 28, 1186-1228.

\item[]  F\`eve, P. et A. Guay (2010), ``Identification of Technology Shocks in Structural VARs'',
{\it Economic Journal,} 120, 1284--1318.

%\item[] Ambler S., A. Guay et L. Phaneuf, ''Labor Market
%Frictions and Endogenous Business Cycle Propagation.'' {\it Journal of Monetary Economics.}

\item[] F\`eve, P. et A. Guay (2009), ``The Response of Hours to a Technology Shock: a Two-Step 
Structural VAR Approach'', 
{\it Journal of Money, Credit and Banking}, 41, 987--1013.



\item[] Bec, F,  A. Guay et E. Guerre (2008), ``Adaptive Consistent
Unit Root Tests Based on Autoregressive Threshold Model'', {\it
Journal of Econometrics,} 142,
94-133.



\item[] Dridi R., A. Guay et E. Renault (2007), ``Indirect Inference and
Calibration of Dynamic Stochastic General Equilibrium Models'', {\it
Journal of Econometrics,} 136, 397-430.


\item[] Guay, A. et E. Guerre (2006), ``A Data-Driven Nonparametric
Specification Test for Dynamic Regression Models,'' {\it Econometric
Theory,} 22, 543-586.


\item[] Guay, A. et P. St-Amant (2005), ``Do Hodrick-Prescott and Baxter-King Filters
Provide a Good Approximation of Business Cycles?" {\it Annales
d'\'Economie et de Statistique,} 77, 133--156.


\item[] Fauvel, Y., A. Guay et A. Paquet (2005), ``Les neuf vies de la courbe de Phillips
am\'ericaine: r\'eincarnations ou r\'esilience," {\it L'Actualit\'e
\'economique,} 81, 665--691.


\item[] Ghysels E. et A. Guay (2004), ``Testing for Structural Change in the Presence of
Auxiliary Models," {\it Econometric Theory,} 20, 1168-1202.


\item[] Guay A. (2003), ``Optimal Predictive Tests," {\it
Econometric Reviews,} 22, 379-410.

\item[] Ghysels, E., et A. Guay (2003), ``Structural Change Tests for Simulated
Method of Moments", {\it Journal of Econometrics,} 115, 91-123.

\item[] Guay, A. et O. Scaillet (2003), ``Indirect Inference, Nuisance Parameter and
Threshold Moving Average Models", {\it Journal of Business and Economic Statistics,} 21, 122--132.

\item[] Guay, A., R. Luger et Z. Zhu (2002), ``The New Phillips Curve in Canada,''
dans {\it Price Adjustment and Monetary Policy,} Bank of Canada.

\item[] Dupasquier, C, A. Guay et P. St-Amant (1999), ``A Comparison of Alternative Methodologies for Estimating Potential Output and the Output Gap," {\it Journal of Macroeconomics,}  21, 577-595.

\item[] Ghysels, E., A. Guay et A. Hall (1998), ``Predictive Tests for Structural
Change with Unknown Breakpoint", {\it Journal of Econometrics,} 82, 209-233.

\item[] Beaudry, P. et A. Guay (1996), ``What Do Interest Rates Reveal about the
Functioning of Real Business Cycle Models ?", {\it Journal of Economic Dynamics and Control,} 20, 1661-1682.

\end{description}

%\centerline{\bf Papier en r\'evision:}
%\vspace{.2 cm}
%
%\begin{description}
%\item "When is Nonfundamentalness in VARs a Real Problem? An Application to News Shocks" (avec Paul Beaudry, Patrick Fève and Frank Portier),
%NBER, Working Papers 21466. {\it European Economic Review.} 
%\end{description}
%
\centerline{\bf Articles en soumission:}
\vspace{.2 cm}
%
%\begin{description}
\begin{description}

\item[] Beaudry, P., F. Collard, P. F\`eve, A. Guay et F. Portier (2026), ``Exploiting Dynamic Identification in VARs'', r\'evision demand\'ee, {\it Journal of Monetary Economics.}



\item[] Guay, A., F. Pelgrin et S. Surprenant (2026), ``Max Share Identification for Structural VARs in Levels: There Is No Free Lunch!'', r\'evision demand\'ee, {\it Journal of Econometrics.}

\item[] Guay, A. et D. Stevanovi\'c (2026), ``A Spectral Framework for Non-Gaussian SVARs'', soumis.

\end{description}

\vspace{.3cm}

\centerline{\bf Travaux en cours:}
\vspace{.2 cm}

\begin{description}

\item[] Guay, A. et D. Stevanovi\'c, ``Identifying Dynamic Factors in Multiway Data: An Orthogonal CP Decomposition''.

\item[] Guay, A., ``Rank Tests and Identification Strength for Moment Operators: A Spectral Perturbation Approach''.

\item[] Guay, A. et D. Stevanovi\'c, ``Non-Gaussian Factor Models via Spectral Decomposition''.


%\item[] Guay, A. and F. Pelgrin (2024), ``Structural Vector Autoregressive and the Asymptotic Least Squares Principle: A General Framework,'' soumis \`a la revue {\it Journal of Econometrics.}

%\item[]  Collard, F., P. F\`eve and A. Guay (2024),  ``  Believe it or not, it’s all about Beliefs!'' soumis \`a la revue {\it American Economic Journal: Macroeconomics}

\end{description}






\vspace{.3cm}


%\newpage

\centerline{\bf Financement de la recherche}
%\begin{description}

\begin{itemize}

\item CRSH-Savoir, projet intitul\'e : Multidimensional analysis in economics: Estimation, Inference and Applications (2026-2029).

\item CRSH-Connexion, avec Dalibor Stevanovi\'c pour l'organisation de la conf\'erence annuelle de ISF (International Institute of Forecasters) Conference, \$25,000 (2026).


\item CRSH-Connexion, avec Dalibor Stevanovi\'c pour l'organisation de la conf\'erence annuelle du NBER-NSF Time Series Conference, \$25,000 (2023).

\item CRSH-Savoir, individuel, subvention pour un montant
de \$92,416 (2018-2023).

\item CRSH-Connexion, avec Dalibor Stevanovi\'c et Benoit Perron, pour l'organisation de la conf\'erence annuelle de l'International Association of Applied Econometrics (IAAE), \$25,000.

\item FQRSC \'equipe avec Hafedh Bouakez (HEC), Michel Normandin (HEC) et Francisco Ruge-Murcia (McGill). Subvention pour quatre ans au montant
individuel de \$14,000 par ann\'ee (2017-2022) pour un montant total de 224,000\$ pour l'\'equipe pour les quatre ann\'ees.

\item FQRSC \'equipe avec Hafedh Bouakez (HEC), Michel Normandin (HEC), Louis Phaneuf (ESG-UQAM) et Frederico Ravenna (HEC). Subvention pour quatre ans au montant
individuel de \$21,000 par ann\'ee (2012-2016).

\item ANR (Agence Nationale de la Recherche, France) \'equipe avec Patrick F\`eve (TSE), Franck Portier (TSE) et Martial Dupaigne (TSE et Montpellier). Subvention pour quatre ans au montant
total de 150,000 euros  (2014-2018).


\item FQRSC \'equipe avec Jean Boivin (HEC), Hafedh Bouakez (HEC), Michel Normandin (HEC) et Louis Phaneuf (ESG-UQAM). Subvention pour quatre ans au montant
individuel de \$16,000 par ann\'ee (2008-2012).

\item CRSH individuel, subvention pour trois ans pour un montant
annuel de \$20,000 (2008-2011).

\end{itemize}


\vspace{.3cm}

\centerline{\bf Direction de th\`eses et m\'emoires}

\vspace{.3cm}

J'ai dirig\'e (ou en cours de direction) :  
\begin{itemize}

\item 67 m\'emoires de ma\^{\i}trise 
\item 10 th\`eses de doctorat
\end{itemize}
%\centerline{\bf Direction de m\'emoires depuis 2008}

%\vspace{.3cm}

%\begin{itemize}

%

%
%\item{\bf Olivier Gervais}, {\em Les changements structurels peuvent-ils \^etre la cause de la grande mod\'eration de l'\'economie am\'ericaine?},
%codirection
%avec Louis Phaneuf, d\'ep\^ot en
%juin 2008.

%
%\item{\bf  \'Eric Morin}, {\em \'Evolution de la part de la valeur
%boursi\`ere des fiducies de revenus attribuable à la prime li\'e \`a
%leur statut juridique: une approche \'econom\'etrique}, codirection
%avec Doug Hodgson,  d\'ep\^ot en
%ao\^ut 2008.

%
%\item{\bf Pierre-Guy Sylvestre}, {\em L'impact de la privatisation et de la r\`eglementation dans l'industrie de l'\'electricit\'e}, codirection avec Nicolas Marceau,  d\'ep\^ot en ao\^ut 2009.

%

%
%\item{\bf Myriam Moisan}, {\em La pr\'evision de l'inflation au Canada}, d\'ep\^ot en f\'evrier 2010.

%\item{\bf Manal Nessim}, {\em La r\`egle deTaylor: varier le taux cible d'inflation}, 
%codirection
%avec Louis Phaneuf, mai 2012.

%\item{\bf Marc Simard}, {\em \'Etude comparative de la croissance \'economique entre la Slovaquie et la R\'epublique th\`eque depuis la r\'evolution de velours }, 
%codirection avec Nicolas Marceau, d\'ep\^ot  en janvier 2012.

%\item{\bf \'Emile Chiasson}, {\em Estimations par des mod\`eles de corr\'elations conditionnellles dynamiques}, 
%d\'ep\^ot en mai 2012.

%\item{\bf Junia Barreau}, {\em L'\'evolution du march\'e du travail au Canada et aux \'Etats-Unis depuis la Grande R\'ecession de 2007--2009},
%codirection
%avec Louis Phaneuf, d\'ep\^ot mai 2014.

%\item{\bf Hamidou Zanr\'e}, {\em Impacts des chocs gouvernementaux sur l'\'economie canadienne}, 
%d\'ep\^ot mars 2015..

%\item{\bf Laurence Savoie-Chabot}, {\em Chocs de sentiments, la confiance des producteurs et les fluctuations macro\'economiques  am\'ericaines}, 
%d\'ep\^ot mars 2015.

%\item{\bf Jocelyn Bissonnette}, {\em Structure \`a terme des taux d'int\'er\^et et ''news shocks''}, 
%travail en cours.

%\item{\bf Isabelle Boileau}, {\em Chocs de sentiments et les fluctuations macro\'economiques canadiennes}, 
%travail en cours.

%\item{\bf Samuel Brien}, {\em L'importance des chocs d'incertitude pour les fluctuations macro-\'economiques am\'ericaines}, 
%travail en cours.

%\item{\bf Vincent Beaus\'ejour}, {\em L'impact des fluctuations \'economiques sur les in\'egalit\'es de revenu aux \'Etats-Unis }, 
%travail en cours.

%\item{\bf Beno\^it Vincent}, {\em Indicateur de l'activit\'e r\'eelle au Canada}, codirection avec Dalibor Stevanovic,
%travail en cours.

%\item{\bf Martinien Dansour}, {\em Un mod\`ele \`a facteurs pour le Canada}, codirection avec Dalibor Stevanovic,
%travail en cours.

%\item{\bf Olivier Fortin-Gagnon}, {\em Mesures de l'incertitude pour le Canada}, codirection avec Dalibor Stevanovic,
%travail en cours.

%\end{itemize}

%\vspace{.3cm}

%\centerline{\bf Direction de th\`eses depuis 2008}

%\vspace{.3cm}

%\begin{itemize}

%

%
%\item{\bf Pierre Chauss\'e}, {\em M\'ethode de vraisemblance empirique et continuum de moments}, d\'ep\^ot en f\'evrier 2012.

%
%\item{\bf J\'er\'emy Chaudourne}, {\em Trois essais en \'econom\'etrie }, d\'ep\^ot septembre 2013.

%\item{\bf Jean-Gardy Victor}, {\em Essais sur la dynamique du cycle \'economique, les frictions financi\`eres et la politique \'economique }, travail en cours.

%
%\end{itemize}


\vspace{.3cm}


\centerline{\bf Responsabilit\'es administratives}
\begin{description}

\vspace{.3 cm}

\item Cotitulaire de la Chaire en macro\'economie et pr\'evisions de l'ESG-UQAM (2020--). Cette Chaire, qui a d\'ebut\'e ses activit\'es au d\'ebut de 2020, est financ\'ee par le Minist\`ere des Finances et le Minist\`ere de l'Emploi et de la Solidarit\'e sociale du Gouvernement du Qu\'ebec, Hydro-Qu\'ebec, la Caisse de d\'ep\^ot et placement du Qu\'ebec et de l'Autorit\'e des march\'es financiers .

\item Directeur du D\'epartement des sciences \'economiques (2015--2020), \'Ecole des sciences de la gestion, UQAM.

\item Codirecteur du centre de recherche CIRP\'EE et directeur de
la composante CIRP\'EE-UQAM, 2005 \`a 2010. Ce centre de recherche interuniversitaire
\'etait financ\'e par le FQRSC en tant que regroupement strat\'egique.


\item Membre du comit\'e de la recherche de l'ESG (
2005-2006, 2006-2007, 2007-2008, 2008-2009, 2009-2010).

\item Membre du comit\'e des partenaires du Consortium sur l'innovation,
les performances et le bien-\^etre dans l'\'economie du savoir
(CIBL'es) (2005-2006, 2006-2007, 2007-2008 et 2008-2009).

\item Membre du comit\'e ex\'ecutif du d\'epartement (2003-2004, 2005-2006, 2008-2009, 2009-2010, 2011-2012, 2012-2013, 2015-2016, 2016-2017, 2017-2018, 2018-2019, 2019-2020, 2020-2021, 2021-2022, 2022-2023, 2023-2024, 2024-2025).

\item Membre du comit\'e d'examen des candidatures professorales (2003-2004, 2005-2006,2008-2009, 2011-2012, 2012-2013 et 2013-2014, 2015-2016, 2016-2017, 2017-2018, 2019-2020, 2021-2022, 2024-2025).

\item Membre du comit\'e d'\'evaluation des professeurs du
d\'epartement. (2004-2005, 2006-2007, 2012-2013 (substitut), 2013-2014 (substitut), 2014-2015, 2015-2016, 2016-2017, 2017-2018).

\item Membre du comit\'e des \'etudes de cycles sup\'erieurs du d\'epartement (2003-2004, 2004-2005, 2005-2006, 2006-2007, 2007-2008, 2008-2009 et 2009-2010 ).

\item Membre du comit\'e de recherche et d'orientation du
d\'epartement (2005-2006, 2006-2007, 2007-2008, 2008-2009, 2009-2010, 2015-2016, 2016-2017, 2017-2018, 2018-2019, 2019-2020, 2020-2021, 2021-2022, 2022-2023, 2024-2025).

\item Responsable du placement des \'etudiants des cycles
sup\'erieurs du d\'epartement (2003-2004, 2004-2005, 2005-2006, 2006-2007, 2007-2008, 2008-2009, 2009-2010 et 2011-2012).

\item Repr\'esentant du d\'epartement \`a la ma\^itrise en finance appliqu\'ee et DESS en finance (2009-2010).

\item Conseil d'administration de l'Institut des sciences de
l'environnement (2015-2016, 2016-2017, 2017-2018).

\item Membre du comit\'e sectoriel institutionnel d'octroi de bourses CRSH pour la ma\^{\i}trise
et le doctorat \`a l'UQAM (2004).


\item Membre du protocole CSN/FTQ/UQAM (2004-2005, 2005-2006,
2006-2007 et 2007-2008).

\item Membre du groupe de discussion conjoint CSN/FTQ/UQAM  portant sur des questions \'economiques (2004-2005, 2005-2006 et
2006-2007).

\item Membre du conseil syndical (SPUQ) (2007-2008 et 2008-2009, 2009-2010, 2011-2012, 2012-2013, 2013-2014 et 2014-2015)

\item Membre du comit\'e de mentorat du d\'epartement (2013-2014 et 2014-2015, 2020-2021, 2021-2022, 2022-2023, 2023-2024,2024-2025)
\end{description}



\vspace{.3cm}




\vspace{1cm} \centerline{\bf Autres}

\vspace{.3cm}

\begin{description}
\item Co\'editeur d'un num\'ero sp\'ecial du {\it Journal of Econometrics} portant sur l'estimation \`a haute fr\'equence \`a para\^{\i}tre en 2026.
\item Organisateur (avec Dalibor Stevanovic) de la conf\'erence \`a venir de ISF (International Institute of Forecasters)  Conference, octobre 2026.
\item Organisateur (avec  Marine Carrasco et  Dalibor Stevanovic) de la conf\'erence en l'honneur du professeur  Eric Ghysels, 10-11 Mai 2024.
\item Organisateur (avec Dalibor Stevanovic) de la conf\'erence du NBER-NSF Time Series Conference, octobre 2023.
\item Organisateur local (avec Dalibor Stevanovic) de la conf\'erence du North American Summer Meeting de l'Econometric Society, 10 au 13 juin 2021.
\item Organisateur local (avec Benoit Perron et Dalibor Stevanovic) de la conf\'erence annuelle de l'International Association of Applied Econometrics (IAAE), 26 au 29 juin 2018.
\item Pr\'esident sortant de la Soci\'et\'e canadienne de science 
\'economique (SCSE, 2016-2017).
\item Pr\'esident de la Soci\'et\'e canadienne de science 
\'economique (SCSE, 2015-2016).
\item Pr\'esident d\'esign\'e de la Soci\'et\'e canadienne de science 
\'economique (SCSE, 2014-2015) et organisateur du congr\`es annuel, UQAM, 2015.
\item Vice-pr\'esident de la Soci\'et\'e canadienne de science
\'economique (SCSE) (2004-2005, 2005-2006, 2006-2007, 2007-2008 et 2008-2009).
\item Membre du comit\'e d'octroi de subventions FQRSC aux jeunes chercheurs (2004 et 2005).
\item Pr\'esident du comit\'e d'octroi de subventions FQRSC aux jeunes chercheurs (2009, 2010 et 2011).
\item Membre du comit\'e d'organisation de la conf\'erence mondiale 2015 
de l'Econometric Society.
\item Membre du comit\'e d'organisation de la conf\'erence annuelle
du Canadian Econometric Study Group en 2006.
\item Consultant pour KPMG (2015-2018).
\item Consultant pour la Caisse de d\'ep\^ot et placement du Qu\'ebec (2014): projet portant sur la pr\'evision de l'inflation am\'ericaine \`a l'aide de donn\'ees \`a diff\'erentes fr\'equences.
\item Consultant pour la Banque du Canada dans le cadre de deux projets intitul\'es: 
      \begin{itemize}
            \item {\em Another Look at Housing Wealth Effects in Canada.} 
            \item {\em Impact des prix des maisons sur la consommation des biens non durables des m\'enages.} 
      \end{itemize}
\item Consultant pour l'OCDE. Deux documents de travail \`a para\^{\i}tre:
      \begin{itemize}
            \item {\em HP filtering and SVAR models: An assessment,} avec F. Pelgrin, 75 pages.
            \item {\em SVAR models and the short-run resilience effect,} avec F. Pelgrin, 82 pages.
      \end{itemize}

\item Textes de vulgarisations:
\begin{itemize}
            \item Quotidiens
                    \begin{enumerate}
                           \item {\em Pour que le rem\`ede n'ach\`eve pas le patient,} avec Marie Connally, Catherine Haeck, Charles S\'eguin, Dalibor Stevanovic, Marie-Louise Leroux et Ariana Degan, Le Devoir, d\'ecembre 2014.
                           \item {\em Le budget de l'an 1 de Fran\c cois Legault est un document \'equilibr\'e et utile,} avec Nicolas Marceau dans La Presse, 13 mai 2005.            
                           \item {\em Depuis 25 ans, le mod\`ele qu\'eb\'ecois a permis \`a notre \'economie de combler une portion importante de son retard,} avec Nicolas Marceau dans La Presse, 13 d\'ecembre 2004.
                     \end{enumerate}      
            \item Monographie et chapitres de livre:
                    \begin{enumerate}
                            \item {\em Le Qu\'ebec n'est pas le cancre \'economique qu'on dit,} avec Nicolas Marceau, dans L'annuaire du Qu\'ebec 2005, Montr\'eal, \'Editions Fides, 66--83, 2004.
                   \end{enumerate}                  
            \end{itemize}                  
\end{description}

\vfill\eject
\end{document}
